\documentclass[11pt,a4paper]{moderncv}
\moderncvstyle{classic}
\moderncvcolor{black}
\definecolor{firstnamecolor}{rgb}{0,0,0}

\usepackage[utf8]{inputenc}
\usepackage[T1]{fontenc}
\usepackage[Hscale=0.7, top=2cm, bottom=2cm]{geometry}
\usepackage{microtype}
\usepackage{hyperref}

\renewcommand{\rmdefault}{ppl}
\renewcommand{\sfdefault}{phv}
\renewcommand*{\namefont}{\fontsize{26}{18}\mdseries}
\microtypesetup{expansion=false}

\input{../../secrets/personal_info_dk.tex}

\begin{document}
\makecvtitle

Dear members of the hiring committee,

I am applying for the PhD position in the AI:HARDWARE Lab because it matches the work I want to do next: rigorous machine learning research under real implementation constraints. My background at KTH combines Engineering Physics with an MSc specialisation in Statistical Learning and Data Analytics, and my degree project on \href{https://kth.diva-portal.org/smash/get/diva2:1431327/FULLTEXT02.pdf}{probabilistic dose-volume histogram and spatial dose prediction} gave me experience of formulating a modelling problem, testing assumptions, and working through evaluation in detail.

At RaySearch, I worked as a researcher on automated radiotherapy planning using deep learning, probabilistic modelling, and multi-objective optimisation. That work included pipelines for latent anatomical features from segmented 3D CT scans, as well as sparse Gaussian-process models for dose prediction. The important part for me was not only the modelling, but the way the work had to be grounded in clinical practice. Collaborating with medical physicists and clinicians forced me to be precise, to justify modelling choices, and to keep the research tied to a real decision process. That is the sort of research setting I am looking for in a PhD.

My more recent industry work strengthens the implementation side of the same profile. At Valcon I architected and implemented an end-to-end production pipeline for a real-time deep neural network, and I built an invoice quality-control system with around two-second response time and 100\% recall in the first two weeks of production. At Signum I currently work on embedding-based similarity search and matching workflows, which keeps me close to applied ML system design. At Valtech, I worked with business stakeholders on forecasting and clustering outputs. Across these roles, I have had to move between model design, implementation, and explanation.

I would bring a strong quantitative base, independent research experience, and practical ML implementation maturity to the PhD environment. That combination is relevant to reliable system work, and I would value the chance to contribute it while developing further within your group.

Thank you for considering my application. I would welcome the opportunity to discuss how my background could contribute to the project.

Yours sincerely,

Ivar Cashin Eriksson

\end{document}
